%%%%%%%%%%%%%%%%%%%%%%%%%%%%%%%%%%%%%%%%%%%%%%%%%%%%%%%%%%%%%%%%%%%%%%%%
% Plantilla TFG/TFM
% Universidad de Murcia. Facultad de Informática
% Realizado por: José Manuel Requena Plens
% Modificado: Pablo José Rocamora Zamora
% Contacto: pablojoserocamora@gmail.com
%%%%%%%%%%%%%%%%%%%%%%%%%%%%%%%%%%%%%%%%%%%%%%%%%%%%%%%%%%%%%%%%%%%%%%%%


% Lista de acrónimos (se ordenan por orden alfabético automáticamente)

% La forma de definir un acrónimo es la siguiente:
% \newacronyn{id}{siglas}{descripción}
% Donde:
% 	'id' es como vas a llamarlo desde el documento.
%	'siglas' son las siglas del acrónimo.
%	'descripción' es el texto que representan las siglas.
%
% Para usarlo en el documento tienes 4 formas:
% \gls{id} - Añade el acrónimo en su forma larga y con las siglas si es la primera vez que se utiliza, el resto de veces solo añade las siglas. (No utilices este en títulos de capítulos o secciones).
% \glsentryshort{id} - Añade solo las siglas de la id
% \glsentrylong{id} - Añade solo la descripción de la id
% \glsentryfull{id} - Añade tanto  la descripción como las siglas

\newacronym{api}{API}{Application Programming Interface}
\newacronym{kg}{KG}{Knowledge Graph}
\newacronym{llm}{LLM}{Large Language Model}
\newacronym{json}{JSON}{JavaScript Object Notation}
\newacronym{nt}{NT}{N-Triples}
\newacronym{nlp}{NLP}{Natural Language Processing}
\newacronym{pln}{PLN}{Procesamiento de Lenguaje Natural}
\newacronym{rag}{RAG}{Retrieval-Augmented Generation}
\newacronym{rdf}{RDF}{Resource Description Framework}
\newacronym{sparql}{SPARQL}{Simple Protocol and RDF Query Language}
\newacronym{ui}{UI}{User Interface}
